%---------- Inleiding ---------------------------------------------------------

% TODO: Is dit voorstel gebaseerd op een paper van Research Methods die je
% vorig jaar hebt ingediend? Heb je daarbij eventueel samengewerkt met een
% andere student?
% Zo ja, haal dan de tekst hieronder uit commentaar en pas aan.

\paragraph{Opmerking}

Dit voorstel is gebaseerd op het onderzoeksvoorstel dat werd geschreven in het
kader van het vak Research Methods dat ik vorig academiejaar heb
uitgewerkt met medesturent Yu-lan Heremans als mede-auteur.
 

\section{Inleiding}%
\label{sec:inleiding}

\textcolor{blue}{Het correct inschatten van het nodige personeel is een cruciale uitdaging voor kleine detailhandelaars zoals Bakkerij Bossu in Nazareth. Onvoldoende personeel kan leiden tot overbelasting van het team en tot lange wachttijden, die al snel invloed kunnen hebben op de klanttevredenheid, terwijl te veel personeel onnodige kosten veroorzaakt. Voor bakkers is het dus essentieel om de personeelsplanning efficiënt en nauwkeurig toe te passen.}

Om dit probleem aan te pakken, focust deze bachelorproef zich op het gebruik van operationele en contextuele data in combinatie met mach\-ine-learningmodellen om de personeelsbehoefte van een bakkerij aan te tonen op basis van voorspelde verkoopcijfers. De kern van dit onderzoek is dat nauwkeurige verkoopvoorspellingen de bakker in staat stellen om optimaal personeel in te plannen, waardoor zowel de kosten als de klanttevredenheid gemaximaliseerd worden.

De onderzoeksvraag luidt: \textcolor{blue}{ “Hoe kunnen operationele en contextuele data van bakkers in combinatie met machine-learningmodellen worden ingezet om de benodigde hoeveelheid personeel te bepalen op basis van verkoopvoorspellingen?”}

Om deze centrale vraag correct te kunnen beantwoorden, worden enkele deelvragen onderzocht. Om de datastructuur te begrijpen, is het interessant om de deelvraag: \textcolor{blue}{ “Hoe sterk variëren de verkoopvolumes tijdens weekends en feestdagen?”} te onderzoeken. Een andere belangrijke deelvraag richt zich op de keuze van het AI-model: \textcolor{blue}{welk machine-learningmodel kan de verkoop het nauwkeurigst voorspellen en daarmee de planning optimaliseren?} Ook de keuze van de te gebruiken data is essentieel voor de voorspelling. \textcolor{blue}{Welke operationele en contextuele data zijn relevant om betrouwbare voorspellingen te maken?} Deze deelvraag houdt hier rekening mee.

De doelstelling van deze bachelorproef is het ontwikkelen van een proof-of-concept voorspellingsmodel dat Bakkerij Bossu ondersteunt bij het plannen van personeel op basis van verkoopvoorspellingen. Daarnaast wordt een kort adviesrapport opgesteld met aanbevelingen voor de planning.


%---------- Stand van zaken ---------------------------------------------------

\section{Literatuurstudie}%
\label{sec:literatuurstudie}

De inzet van artificiële intelligentie en dataengineering in personeelsplanning krijgt steeds meer aandacht, vooral binnen sectoren waar de behoefte aan personeel sterk fluctueert. Deze literatuurstudie bespreekt relevante inzichten en identificeert de belangrijkste uitdagingen, maar ook de opportuniteiten voor het gebruik van AI binnen personeelsplanning, met specifieke aandacht voor de bakkerij-industrie.

\subsection*{1. Vergelijken modelprestaties}
De literatuur toont een duidelijk onderscheid in prestaties tussen traditionele lineaire modellen en moderne ensemble-methoden. \cite{GangulyMukherjee2024} vergeleek verschillende machine-learningmodellen voor retailverkoopvoorspelling en stelde concrete verschillen vast in prestatie: Random Forest behaalde een R² van 0.945 met een RMSE van 282.07, terwijl lineaire regressie slechts een R² van 0.531 bereikte. Deze resultaten tonen aan dat ensemble-methoden aanzienlijk sterker presteren bij het modelleren van verkooppatronen.  

Gradient Boosting behoort ook tot de best presterende modellen. \cite{GangulyMukherjee2024} rapporteert een R² van 0.942, terwijl \cite{Kang2023} aantoont dat Gradient Boosting in hun studie marginaal beter presteerde dan Random Forest. Deze consistente prestaties op verschillende datasets wijzen erop dat ensemble-methoden geschikt zijn voor het voorspellen van retailverkopen.

\subsection*{2. Rol van operationele en contextuele data}
\cite{Geboers2012} onderzocht in zijn case study bij Bakkersland B.V. de impact van contextuele variabelen op verkoopvoorspellingen voor de broodindustrie. Het concludeert dat het toevoegen van contextuele data, zoals weersinvloeden, feestdagen en lokale evenementen, de nauwkeurigheid significant verbetert ten opzichte van modellen die enkel historische verkoopdata gebruiken.

\subsection*{3. Van voorspelling naar personeelsplanning}
\cite{Geboers2012} beschrijft in zijn case study hoe verkoopvoorspellingen kunnen worden vertaald naar productieplanning. Voor personeelsplanning vereist dit het definiëren van capaciteitsregels: hoeveel klanten kan één medewerker per tijdseenheid bedienen, wat is de minimale bezetting ongeacht voorspelde drukte, en welke buffer is nodig voor onzekerheid?

\subsection*{4. Conclusie}
De literatuur toont consistent aan dat ensemble-methoden zoals Random Forest en Gradient Boosting significant beter presteren dan traditionele lineaire modellen voor het voorspellen van verkoop.  

Voor bakkerijen is het integreren van zowel operationele data als contextuele data cruciaal voor een nauwkeurige voorspelling. De combinatie van deze data en machine-learningmodellen stelt kleine ondernemingen in staat om de personeelsplanning te optimaliseren.  

De uitdaging ligt in het praktisch implementeren van deze technologieën. Een proof-of-concept moet niet alleen robuust zijn, maar ook praktisch toepasbaar en gebruiksvriendelijk voor de eindgebruiker.


% Voor literatuurverwijzingen zijn er twee belangrijke commando's:
% \autocite{KEY} => (Auteur, jaartal) Gebruik dit als de naam van de auteur
%   geen onderdeel is van de zin.
% \textcite{KEY} => Auteur (jaartal)  Gebruik dit als de auteursnaam wel een
%   functie heeft in de zin (bv. ``Uit onderzoek door Doll & Hill (1954) bleek
%   ...'')

%---------- Methodologie ------------------------------------------------------
\section{Methodologie}%
\label{sec:methodologie}

Het doel van deze bachelorproef is het ontwikkelen van een proof-of-concept dat bakkerij Bossu ondersteunt bij het plannen van zijn personeel. Om de onderzoeksvraag te beantwoorden, wordt het onderzoek opgedeeld in onderstaande iteratieve fasen, die elkaar overlappen en iteratieve verbeteringen toelaten.

\paragraph{Fase 1: Literatuurstudie}
Deze fase omvat een grondige analyse van literatuur over machine-lea\-rningmodellen voor verkoopvoorspellingen, personeelsplanning in kleine ondernemingen en het gebruik van operationele en contextuele data. Het doel is het selecteren van geschikte modellen voor de casus van bakkerij Bossu, en er kennis over opdoen om deze correct toe te passen op de verzamelde data. Maar ook vaststellen welke soorten data invloed zullen hebben op de verkoop. De deliverable van deze fase is een overzicht dat de selectiecriteria voor machine-learningmodellen en de relevante data-elementen beschrijft. Dit stadium wordt gepland om twee weken in beslag te nemen en dient als basis voor alle verdere fasen.

\paragraph{Fase 2: Data-verzameling en -voorbereiding}
Deze stap richt zich op het verzamelen en voorbereiden van operationele en contextuele data bij bakkerij Bossu. Deze data wordt verzameld uit de winkel van Bossu. Hiertoe behoren verkoopcijfers, openingstijden, seizoen, feestdagen en eventuele lokale evenementen. De data wordt geanalyseerd om trends en correlaties te identificeren zodat de meest rendabele variabelen overblijven. Deze data-analyse wordt uitgevoerd op een Pandas-dataframe en gevisualiseerd via matplotlib. De overblijvende dataset na deze analyse is de deliverable. Dit staat gepland om twee weken in te nemen.

\paragraph{Fase 3: Experimenten met machine-learningmodellen}
Op basis van de literatuurstudie en data-analyse worden meerdere machine-learnin\-gmodellen geselecteerd waaronder verschillende regressiemodellen en random forests. Deze modellen worden getraind en getest op de verzamelde data om de nauwkeurigheid te kunnen evalueren. Deze evaluatie verloopt aan de hand van objectieve metrics zoals RMSE en MAE. Daarna worden er iteratieve verbeteringen doorgevoerd door de hyperparameters aan te passen en de modellen opnieuw te trainen. De deliverable van deze fase is een rapport waarin de prestaties van de verschillende modellen worden vergeleken en het meest geschikte model voor de casus wordt gekozen. De planning voor deze fase bedraagt drie weken.

\paragraph{Fase 4: Proof-of-concept ontwikkeling}
In fase 4 wordt een proof-of-concept ontwikkeld dat de voorspelde verkoopcijfers omzet naar een berekening van de benodigde personeelsleden. Het prototype richt zich niet op de volledige personeelsplanning, zo is er nog ruimte voor de bakker voor eigen invloeden. Er worden op voorhand capaciteitsregels gedefinieerd, zoals het aantal producten dat één medewerker per tijdseenheid kan verwerken. Op basis van deze regels wordt voorspeld hoeveel werknemers nodig zullen zijn. Dit kan de bakker dan meenemen in zijn planning en zelf werknemers inplannen voor deze hoeveelheden. Het werkende proof-of-concept is de deliverable voor deze fase en deze zal drie weken in beslag nemen.


%---------- Verwachte resultaten ----------------------------------------------
\section{Verwacht resultaat, conclusie}
\label{sec:verwachte_resultaten}

In dit onderzoek wordt verwacht dat één van de geteste machine-learningmodellen duidelijk de beste prestaties zal leveren voor het voorspellen van de verkoop bij Bakkerij Bossu.
Op basis van de literatuur wordt verwacht dat het voorspellen van de verkopen accuraat kan worden uitgevoerd. Daarnaast suggereren meerdere studies dat lineaire modellen doorgaans minder goed presteren in een retailcontext, terwijl modellen zoals Random Forest en Gradient Boosting vaak tot de best scorende behoren \cite{GangulyMukherjee2024, Kang2023}. Daarom wordt verwacht dat ook in deze casus een ensemble-model de hoogste nauwkeurigheid zal behalen.

Hoewel deze verwachtingen gebaseerd zijn op bestaande literatuur, blijft het mogelijk dat de resultaten afwijken. Indien een model minder goed presteert dan verwacht, zal worden onderzocht welke factoren bijdragen aan deze overschatting.